\answer{ex:05:1}{$\ell\approx17\um$; $\approx100$~days: broadcasting across the brain is done by the branching wire, and diffusion only blurs each release site over $\sim20\um$.}
\answer{ex:05:2}{$r=ab/(1+G)$, $S=1/(1+G)$ exactly; $+1.7\,\%$ instead of $+20\,\%$.}
\answer{ex:05:3}{$T^*=1.21\,\text{s}$ for $G=2$ and $0.19\,\text{s}$ for $G=9$: realistically $G\sim2$--$4$, a suppression by a factor of $3$--$5$.}
\answer{ex:05:4}{$\mathrm{CV}=1/\sqrt{2\lambda\tau_c}\approx9\cdot10^{-4}$ for a total rate $\lambda$; a single neuron needs $\tau_c=12.5\,\text{s}$.}
\answer{ex:05:5}{$c\approx0.40$ in both cases; $\mathrm{CV}=0.050$ versus $0.158$, a factor of $\sqrt{10}$. The rates of individual neurons stay proportional to $a_i$: only the sum is regulated.}
\answer{ex:05:6}{$N_1=\overline{A\vee B}$, $N_2=\overline{A\vee N_1}$, $N_3=\overline{B\vee N_1}$, $N_4=\overline{N_2\vee N_3}$, XOR $=N_5=\overline{N_4}$; each switch takes $\tau_c\ln2$, and a path of $4$ gates takes $2.8\,\text{s}$ per XOR.}
\answer{ex:05:7}{(a)~$\tau_\gamma=6/\ln3\approx5.5\,\text{s}$. (b)~$\bar D<4\,\text{s}$ versus $D<2.2\,\text{s}$: an unpredictable delay makes the animal more patient.}
\answer{ex:05:8}{$k_2=1/\tau_d=10\,\text{s}^{-1}$, $k_1=1/\tau_d-1/\tau_\gamma=9.5\,\text{s}^{-1}$; $\tau_\gamma=1/(k_2-mk_1)$: doubling at $+2.6\,\%$ (and at $+5.3\,\%$ the horizon is infinite).}
